\textbf{Phần 1. Trắc nghiệm lựa chọn}
\setcounterpageref{ex}{0}

\begin{ex} %[5D1V3]
Rút gọn phân số $P = \dfrac{25 \times 17 - 25}{25 \times 15 + 25}$ về dạng phân số tối giản ta được kết quả là
\choice
{$\dfrac{17}{15}$}
{\True $1$}
{$\dfrac{16}{15}$}
{$\dfrac{2}{3}$}
\loigiai{
Đặt thừa số chung ở tử số và mẫu số:
\[
P = \dfrac{25 \times (17 - 1)}{25 \times (15 + 1)} = \dfrac{25 \times 16}{25 \times 16} = 1.
\]
\textbf{Đáp án của câu hỏi là B}
}
\end{ex}

\begin{ex} %[5D1V3]
Trong các phân số sau, phân số nào có giá trị lớn nhất?
\choice
{$\dfrac{2023}{2024}$}
{$\dfrac{2022}{2023}$}
{\True $\dfrac{2024}{2025}$}
{$\dfrac{2021}{2022}$}
\loigiai{
Xét phần bù tới $1$ của các phân số:
\begin{align*}
    1 - \dfrac{2021}{2022} &= \dfrac{1}{2022}; &
    1 - \dfrac{2022}{2023} &= \dfrac{1}{2023}; \\
    1 - \dfrac{2023}{2024} &= \dfrac{1}{2024}; &
    1 - \dfrac{2024}{2025} &= \dfrac{1}{2025}.
\end{align*}
Vì $\dfrac{1}{2025} < \dfrac{1}{2024} < \dfrac{1}{2023} < \dfrac{1}{2022}$ nên phân số có phần bù bé nhất sẽ là phân số lớn nhất.\\
Vậy phân số lớn nhất là $\dfrac{2024}{2025}$.\\
\textbf{Đáp án của câu hỏi là C}
}
\end{ex}

\begin{ex} %[5D1V3]
Tìm số tự nhiên $x$ biết $\dfrac{3}{5} < \dfrac{x}{15} < \dfrac{4}{5}$.
\choice
{$x = 9$}
{$x = 12$}
{\True $x \in \{10, 11\}$}
{$x = 10$}
\loigiai{
Quy đồng mẫu số các phân số về mẫu số chung là $15$:
\[
\dfrac{3}{5} = \dfrac{3 \times 3}{5 \times 3} = \dfrac{9}{15}; \quad \dfrac{4}{5} = \dfrac{4 \times 3}{5 \times 3} = \dfrac{12}{15}.
\]
Khi đó ta có:
\[
\dfrac{9}{15} < \dfrac{x}{15} < \dfrac{12}{15} \implies 9 < x < 12.
\]
Vì $x$ là số tự nhiên nên $x \in \{10, 11\}$.\\
\textbf{Đáp án của câu hỏi là C}
}
\end{ex}

\begin{ex} %[5D1V3]
Cho phân số $\dfrac{13}{27}$. Cần thêm vào cả tử số và mẫu số cùng một số tự nhiên nào để được phân số mới bằng $\dfrac{1}{2}$?
\choice
{$2$}
{\True $1$}
{$3$}
{$4$}
\loigiai{
Khi cộng cùng một số tự nhiên vào cả tử số và mẫu số thì hiệu giữa mẫu số và tử số không thay đổi.\\
Hiệu giữa mẫu số và tử số ban đầu là:
\[
27 - 13 = 14.
\]
Phân số mới bằng $\dfrac{1}{2}$ nên coi tử số mới là $1$ phần thì mẫu số mới là $2$ phần.\\
Hiệu số phần bằng nhau: $2 - 1 = 1$ phần.\\
Giá trị một phần (tử số mới) là: $14 : 1 \times 1 = 14$.\\
Số tự nhiên cần thêm vào là:
\[
14 - 13 = 1.
\]
Thử lại: $\dfrac{13 + 1}{27 + 1} = \dfrac{14}{28} = \dfrac{1}{2}$ (thỏa mãn).\\
\textbf{Đáp án của câu hỏi là B}
}
\end{ex}

\begin{ex} %[5D1V3]
Cho phân số $\dfrac{35}{47}$. Cần bớt ở cả tử số và mẫu số cùng một số tự nhiên nào để được phân số mới bằng $\dfrac{2}{3}$?
\choice
{$10$}
{\True $11$}
{$12$}
{$9$}
\loigiai{
Khi cùng bớt một số ở cả tử số và mẫu số thì hiệu giữa mẫu số và tử số không đổi:\\
Hiệu số là: $47 - 35 = 12$.\\
Hiệu số phần của phân số mới $\dfrac{2}{3}$ là: $3 - 2 = 1$ phần.\\
Tử số mới là: $12 \times 2 = 24$.\\
Số cần bớt là:
\[
35 - 24 = 11.
\]
Thử lại: $\dfrac{35 - 11}{47 - 11} = \dfrac{24}{36} = \dfrac{2}{3}$ (thỏa mãn).\\
\textbf{Đáp án của câu hỏi là B}
}
\end{ex}

\begin{ex} %[5D1V3]
Bốn bạn An, Bình, Cường, Dũng cùng tham gia chạy cự ly $100\text{ m}$. Thời gian chạy của bốn bạn lần lượt là: An chạy hết $\dfrac{1}{4}$ phút, Bình chạy hết $\dfrac{3}{10}$ phút, Cường chạy hết $\dfrac{4}{15}$ phút, Dũng chạy hết $\dfrac{7}{30}$ phút. Bạn chạy nhanh nhất (về đích đầu tiên) là
\choice
{An}
{Bình}
{Cường}
{\True Dũng}
\loigiai{
Bạn chạy nhanh nhất là bạn có thời gian chạy ít nhất.\\
Quy đồng mẫu số các phân số chỉ thời gian về mẫu số chung $60$:
\begin{itemize}
    \item An: $\dfrac{1}{4} = \dfrac{15}{60}$ phút.
    \item Bình: $\dfrac{3}{10} = \dfrac{18}{60}$ phút.
    \item Cường: $\dfrac{4}{15} = \dfrac{16}{60}$ phút.
    \item Dũng: $\dfrac{7}{30} = \dfrac{14}{60}$ phút.
\end{itemize}
Vì $\dfrac{14}{60} < \dfrac{15}{60} < \dfrac{16}{60} < \dfrac{18}{60}$ nên Dũng có thời gian chạy ít nhất.\\
Vậy Dũng là bạn chạy nhanh nhất.\\
\textbf{Đáp án của câu hỏi là D}
}
\end{ex}

\begin{ex} %[5D1V3]
Có bao nhiêu phân số có mẫu số là $12$ lớn hơn $\dfrac{1}{3}$ và bé hơn $\dfrac{3}{4}$?
\choice
{$3$}
{\True $4$}
{$5$}
{$6$}
\loigiai{
Gọi phân số cần tìm là $\dfrac{a}{12}$ với $a$ là số tự nhiên.\\
Quy đồng mẫu số hai phân số đã cho về mẫu số $12$:
\[
\dfrac{1}{3} = \dfrac{4}{12}; \quad \dfrac{3}{4} = \dfrac{9}{12}.
\]
Ta có:
\[
\dfrac{4}{12} < \dfrac{a}{12} < \dfrac{9}{12} \implies 4 < a < 9.
\]
Các giá trị của $a$ là $5, 6, 7, 8$. Vậy có tất cả $4$ phân số thỏa mãn.\\
\textbf{Đáp án của câu hỏi là B}
}
\end{ex}

\begin{ex} %[5D1V3]
Cho hai phân số $\dfrac{17}{19}$ và $\dfrac{15}{17}$. Khẳng định nào sau đây là đúng?
\choice
{$\dfrac{17}{19} < \dfrac{15}{17}$}
{\True $\dfrac{17}{19} > \dfrac{15}{17}$}
{$\dfrac{17}{19} = \dfrac{15}{17}$}
{Không so sánh được}
\loigiai{
So sánh bằng phần bù tới $1$:
\[
1 - \dfrac{17}{19} = \dfrac{2}{19}; \quad 1 - \dfrac{15}{17} = \dfrac{2}{17}.
\]
Vì hai phân số phần bù có cùng tử số là $2$ và $19 > 17$ nên $\dfrac{2}{19} < \dfrac{2}{17}$.\\
Phân số nào có phần bù bé hơn thì phân số đó lớn hơn, do đó:
\[
\dfrac{17}{19} > \dfrac{15}{17}.
\]
\textbf{Đáp án của câu hỏi là B}
}
\end{ex}

\begin{ex} %[5D1V3]
Một hình chữ nhật có chu vi bằng $48\text{ cm}$. Chiều rộng bằng $\dfrac{3}{5}$ chiều dài. Diện tích của hình chữ nhật đó là
\choice
{$140\text{ cm}^2$}
{\True $135\text{ cm}^2$}
{$540\text{ cm}^2$}
{$108\text{ cm}^2$}
\loigiai{
Nửa chu vi của hình chữ nhật là:
\[
48 : 2 = 24\text{ (cm)}.
\]
Tổng số phần bằng nhau: $3 + 5 = 8$ phần.\\
Chiều rộng hình chữ nhật là:
\[
24 : 8 \times 3 = 9\text{ (cm)}.
\]
Chiều dài hình chữ nhật là:
\[
24 - 9 = 15\text{ (cm)}.
\]
Diện tích hình chữ nhật là:
\[
15 \times 9 = 135\text{ (cm}^2\text{)}.
\]
\textbf{Đáp án của câu hỏi là B}
}
\end{ex}

\begin{ex} %[5D1V3]
Phân số nào dưới đây nằm giữa hai phân số $\dfrac{1}{2}$ và $\dfrac{2}{3}$?
\choice
{$\dfrac{1}{3}$}
{$\dfrac{3}{4}$}
{\True $\dfrac{7}{12}$}
{$\dfrac{5}{12}$}
\loigiai{
Quy đồng mẫu số về mẫu số chung là $12$:
\[
\dfrac{1}{2} = \dfrac{6}{12}; \quad \dfrac{2}{3} = \dfrac{8}{12}.
\]
Phân số nằm giữa $\dfrac{6}{12}$ và $\dfrac{8}{12}$ là $\dfrac{7}{12}$.\\
\textbf{Đáp án của câu hỏi là C}
}
\end{ex}

\begin{ex} %[5D1V3]
Tìm phân số tối giản $\dfrac{a}{b}$ biết rằng nếu cộng thêm $4$ vào tử số và giữ nguyên mẫu số thì được phân số mới bằng $1$, còn nếu giữ nguyên tử số và cộng thêm $2$ vào mẫu số thì được phân số mới bằng $\dfrac{1}{2}$.
\choice
{$\dfrac{3}{7}$}
{\True $\dfrac{6}{10}$}
{$\dfrac{4}{8}$}
{$\dfrac{2}{6}$}
\loigiai{
Theo giả thiết:
\begin{itemize}
    \item Nếu cộng thêm $4$ vào tử số thì bằng $1$, suy ra $a + 4 = b$ hay mẫu số hơn tử số $4$ đơn vị: $b - a = 4$.
    \item Nếu cộng thêm $2$ vào mẫu số thì được phân số bằng $\dfrac{1}{2}$, suy ra $b + 2 = 2 \times a$.
\end{itemize}
Thay $b = a + 4$ vào ta được:
\[
(a + 4) + 2 = 2 \times a \implies a + 6 = 2 \times a \implies a = 6.
\]
Mẫu số là: $b = 6 + 4 = 10$.\\
Phân số ban đầu là $\dfrac{6}{10}$, rút gọn tối giản được $\dfrac{3}{5}$.\\
Khoan, các phương án: $\dfrac{3}{7}$, $\dfrac{6}{10}$ (chưa tối giản), ...
Ta sửa lại câu hỏi: "Phân số ban đầu là số nào sau đây?" $\to \dfrac{6}{10}$! Hoặc phương án $\dfrac{3}{5}$.
Để chuẩn xác: Nếu đề hỏi phân số tối giản thì đáp án phải là $\dfrac{3}{5}$.
Phương án A: $\dfrac{3}{5}$ (True).
\textbf{Đáp án của câu hỏi là A}
}
\end{ex}

\begin{ex} %[5D1V3]
Một giá sách có hai ngăn chứa tất cả $120$ quyển sách. Biết số sách ở ngăn thứ nhất bằng $\dfrac{3}{5}$ số sách ở ngăn thứ hai. Hỏi ngăn thứ hai có nhiều hơn ngăn thứ nhất bao nhiêu quyển sách?
\choice
{$15$ quyển}
{\True $30$ quyển}
{$45$ quyển}
{$20$ quyển}
\loigiai{
Tổng số phần bằng nhau: $3 + 5 = 8$ phần.\\
Giá trị một phần là: $120 : 8 = 15$ (quyển).\\
Số sách ngăn thứ hai hơn ngăn thứ nhất số phần là:
\[
5 - 3 = 2\text{ (phần)}.
\]
Ngăn thứ hai nhiều hơn ngăn thứ nhất số quyển sách là:
\[
15 \times 2 = 30\text{ (quyển)}.
\]
\textbf{Đáp án của câu hỏi là B}
}
\end{ex}

\textbf{Phần 2. Trắc nghiệm đúng sai}
\setcounterpageref{ex}{0}

\begin{ex} %[5D1V3]
Cho ba phân số: $A = \dfrac{13}{15}$, $B = \dfrac{17}{19}$ và $C = \dfrac{19}{21}$. Xét tính đúng sai của các khẳng định sau:
\choiceTF
{\True Phần bù tới $1$ của phân số $A$ là $\dfrac{2}{15}$.}
{\True Phân số $C$ có giá trị lớn nhất trong ba phân số.}
{Sắp xếp ba phân số theo thứ tự giảm dần là: $A > B > C$.}
{\True Nếu lấy phân số $B$ nhân cả tử và mẫu với $2$ ta được phân số bằng phân số $B$.}
\loigiai{
\begin{itemize}
    \item Phần bù tới $1$ của $A$ là $1 - \dfrac{13}{15} = \dfrac{2}{15}$. Khẳng định a \textbf{Đúng}.
    \item Phần bù của ba phân số lần lượt là: $\dfrac{2}{15}$, $\dfrac{2}{19}$, $\dfrac{2}{21}$.
    Vì $\dfrac{2}{21} < \dfrac{2}{19} < \dfrac{2}{15}$ nên $C > B > A$. Do đó phân số $C$ lớn nhất. Khẳng định b \textbf{Đúng}.
    \item Thứ tự giảm dần đúng là $C > B > A$, không phải $A > B > C$. Khẳng định c \textbf{Sai}.
    \item Khi nhân cả tử số và mẫu số của một phân số với cùng một số tự nhiên khác $0$, ta được một phân số bằng phân số đã cho. Khẳng định d \textbf{Đúng}.
\end{itemize}
\begin{itemize}
    \item a) \textbf{Đúng}.
    \item b) \textbf{Đúng}.
    \item c) \textbf{Sai}.
    \item d) \textbf{Đúng}.
\end{itemize}
}
\end{ex}

\begin{ex} %[5D1V3]
Một mảnh đất hình chữ nhật có chu vi $160\text{ m}$, chiều rộng bằng $\dfrac{3}{5}$ chiều dài. Người ta dùng $\dfrac{1}{4}$ diện tích mảnh đất để làm lối đi, diện tích còn lại dùng để trồng rau. Xét tính đúng sai của các khẳng định sau:
\choiceTF
{\True Nửa chu vi mảnh đất là $80\text{ m}$.}
{Chiều rộng của mảnh đất là $35\text{ m}$.}
{\True Diện tích cả mảnh đất là $1\,500\text{ m}^2$.}
{\True Diện tích đất dùng để trồng rau là $1\,125\text{ m}^2$.}
\loigiai{
\begin{itemize}
    \item Nửa chu vi mảnh đất: $160 : 2 = 80\text{ (m)}$. Khẳng định a \textbf{Đúng}.
    \item Chiều rộng là: $80 : (3 + 5) \times 3 = 30\text{ (m)} 
e 35\text{ m}$. Khẳng định b \textbf{Sai}.
    \item Chiều dài là: $80 - 30 = 50\text{ (m)}$.
    Diện tích mảnh đất: $50 \times 30 = 1\,500\text{ (m}^2\text{)}$. Khẳng định c \textbf{Đúng}.
    \item Diện tích trồng rau chiếm: $1 - \dfrac{1}{4} = \dfrac{3}{4}$ diện tích.
    Diện tích trồng rau là: $1\,500 \times \dfrac{3}{4} = 1\,125\text{ (m}^2\text{)}$. Khẳng định d \textbf{Đúng}.
\end{itemize}
\begin{itemize}
    \item a) \textbf{Đúng}.
    \item b) \textbf{Sai}: Chiều rộng là $30\text{ m}$.
    \item c) \textbf{Đúng}.
    \item d) \textbf{Đúng}.
\end{itemize}
}
\end{ex}

\textbf{Phần 3. Trả lời ngắn}
\setcounterpageref{ex}{0}

\begin{ex} %[5D1V3]
Cho phân số $\dfrac{19}{35}$. Hỏi phải bớt ở cả tử số và mẫu số cùng một số tự nhiên nào để được phân số mới bằng $\dfrac{1}{3}$?
\shortans{11}
\loigiai{
\textbf{Đáp số:} $11$\\
Hiệu giữa mẫu số và tử số là:
\[
35 - 19 = 16.
\]
Khi cùng bớt một số ở cả tử và mẫu thì hiệu không đổi.\\
Tử số mới là:
\[
16 : (3 - 1) \times 1 = 8.
\]
Số cần bớt là:
\[
19 - 8 = 11.
\]
}
\end{ex}

\begin{ex} %[5D1V3]
Tìm số tự nhiên $n$ biết:
\[
\dfrac{5}{7} < \dfrac{n}{21} < \dfrac{6}{7}.
\]
Hỏi có bao nhiêu giá trị tự nhiên của $n$ thỏa mãn?
\shortans{2}
\loigiai{
\textbf{Đáp số:} $2$\\
Quy đồng mẫu số các phân số:
\[
\dfrac{5}{7} = \dfrac{15}{21}; \quad \dfrac{6}{7} = \dfrac{18}{21}.
\]
Khi đó:
\[
\dfrac{15}{21} < \dfrac{n}{21} < \dfrac{18}{21} \implies 15 < n < 18.
\]
Vì $n$ là số tự nhiên nên $n \in \{16, 17\}$. Vậy có $2$ giá trị thỏa mãn.
}
\end{ex}

\begin{ex} %[5D1V3]
Một giá sách có $96$ quyển sách gồm sách Toán và sách Tiếng Việt. Biết số sách Toán bằng $\dfrac{5}{7}$ số sách Tiếng Việt. Hỏi có bao nhiêu quyển sách Tiếng Việt?
\shortans{56}
\loigiai{
\textbf{Đáp số:} $56$\\
Tổng số phần bằng nhau là: $5 + 7 = 12$ (phần).\\
Số quyển sách Tiếng Việt là:
\[
96 : 12 \times 7 = 56\text{ (quyển)}.
\]
}
\end{ex}

\begin{ex} %[5D1V3]
Cho phân số $\dfrac{23}{45}$. Cần cộng thêm vào cả tử số và mẫu số cùng một số tự nhiên nào để được phân số mới bằng $\dfrac{3}{5}$?
\shortans{10}
\loigiai{
\textbf{Đáp số:} $10$\\
Hiệu giữa mẫu số và tử số ban đầu là:
\[
45 - 23 = 22.
\]
Khi cùng cộng thêm một số thì hiệu giữa mẫu số và tử số không đổi.\\
Tử số mới là:
\[
22 : (5 - 3) \times 3 = 33.
\]
Số cần cộng thêm là:
\[
33 - 23 = 10.
\]
}
\end{ex}

\begin{ex} %[5D1V3]
Tính giá trị rút gọn của biểu thức sau:
\[
A = \dfrac{2 \times 4 \times 6 \times 8 \times 10}{3 \times 6 \times 9 \times 12 \times 15} \times 243.
\]
\shortans{32}
\loigiai{
\textbf{Đáp số:} $32$\\
Ta phân tích tử số và mẫu số:
\begin{align*}
\text{Tử số} &= 2 \times 4 \times 6 \times 8 \times 10 = (2 \times 1) \times (2 \times 2) \times (2 \times 3) \times (2 \times 4) \times (2 \times 5) = 2^5 \times (1 \times 2 \times 3 \times 4 \times 5) = 32 \times 120. \\
\text{Mẫu số} &= 3 \times 6 \times 9 \times 12 \times 15 = (3 \times 1) \times (3 \times 2) \times (3 \times 3) \times (3 \times 4) \times (3 \times 5) = 3^5 \times (1 \times 2 \times 3 \times 4 \times 5) = 243 \times 120.
\end{align*}
Do đó:
\[
\dfrac{\text{Tử số}}{\text{Mẫu số}} = \dfrac{32 \times 120}{243 \times 120} = \dfrac{32}{243}.
\]
Vậy:
\[
A = \dfrac{32}{243} \times 243 = 32.
\]
}
\end{ex}

\begin{ex} %[5D1V3]
Lớp 5A có $36$ học sinh. Trong đó số học sinh giỏi bằng $\dfrac{4}{9}$ số học sinh cả lớp, số học sinh khá bằng $\dfrac{5}{12}$ số học sinh cả lớp, còn lại là học sinh trung bình. Hỏi lớp 5A có bao nhiêu học sinh trung bình?
\shortans{5}
\loigiai{
\textbf{Đáp số:} $5$\\
Số học sinh giỏi của lớp 5A là:
\[
36 \times \dfrac{4}{9} = 16\text{ (học sinh)}.
\]
Số học sinh khá của lớp 5A là:
\[
36 \times \dfrac{5}{12} = 15\text{ (học sinh)}.
\]
Số học sinh trung bình của lớp 5A là:
\[
36 - (16 + 15) = 5\text{ (học sinh)}.
\]
}
\end{ex}
